\sectionkicker{Source-backed technical notes}
\subsection*{Detailed Chronology: Technical Notes}
本欄は記事本文で圧縮した一次資料上の情報を、比較・再検証しやすい形で整理したものである。時系列、確認済みの主張、留意点、一次資料URLを掲載する。
\medskip
\subsection*{Theme at a glance}
\addcontentsline{toc}{subsection}{Theme at a glance}
\begin{center}
\footnotesize
\begin{tabularx}{\linewidth}{@{}p{0.20\linewidth}p{0.10\linewidth}p{0.14\linewidth}X@{}}
\toprule
資料 & 位置づけ & 種別 & 時系列 \\
\midrule
DeepSeek-R1 API release & 補足資料 & - & 2025-01-20   \\
OpenAI Operator research preview & 補足資料 & - & 2025-01-23   \\
Qwen2.5-VL open models & 補足資料 & - & 2025-01-26   \\
Qwen2.5-Max API availability & 補足資料 & - & 2025-01-28   \\
OpenAI o3-mini & 補足資料 & - & 2025-01-31   \\
OpenAI deep research & 補足資料 & - & 2025-02-02   \\
Claude 3.7 Sonnet and Claude Code & 補足資料 & - & 2025-02-24   \\
GPT-4.5 research preview & 補足資料 & - & 2025-02-27   \\
Responses API, built-in tools, and Agents SDK & 補足資料 & - & 2025-03-11   \\
Gemma 3 open models & 補足資料 & - & 2025-03-12   \\
Claude web search product preview & 補足資料 & - & 2025-03-20   \\
DeepSeek-V3-0324 API upgrade & 補足資料 & - & 2025-03-24   \\
Gemini 2.5 Pro Experimental & 補足資料 & - & 2025-03-25   \\
GPT-4o native image generation & 補足資料 & - & 2025-03-25   \\
Tracing the thoughts of a large language model & 補足資料 & - & 2025-03-27   \\
Qwen2.5-Omni & 補足資料 & - & 2025-03-27   \\
QVQ-Max visual reasoning & 補足資料 & - & 2025-03-28   \\
Llama 4 Scout and Maverick & 補足資料 & - & 2025-04-05   \\
GPT-4.1 API family & 補足資料 & - & 2025-04-14   \\
GPT-4.5 API deprecation announced & 補足資料 & - & 2025-04-14   \\
OpenAI o3 and o4-mini & 補足資料 & - & 2025-04-16   \\
Llama API limited preview & 補足資料 & - & 2025-04-29   \\
Llama Guard 4 and protection tools & 補足資料 & - & 2025-04-29   \\
Qwen3 & 補足資料 & - & 2025-04-29   \\
Anthropic web search API & 補足資料 & - & 2025-05-07   \\
OpenAI Codex cloud software engineering agent & 補足資料 & - & 2025-05-16   \\
Jules coding agent public beta & 補足資料 & - & 2025-05-20   \\
Veo 3, Imagen 4, Flow, and Lyria 2 & 補足資料 & - & 2025-05-20   \\
Responses API remote MCP and tool expansion & 補足資料 & - & 2025-05-21   \\
Claude Opus 4 and Sonnet 4 & 補足資料 & - & 2025-05-22   \\
Gemini 2.5 Pro and Flash generally available & 補足資料 & - & 2025-06-17   \\
\bottomrule
\end{tabularx}
\end{center}
\normalsize

\begin{technicalnote}{DeepSeek-R1 API release}{補足資料}
\begingroup
% reader-facing Technical Notes break policy
\widowpenalty=10000
\clubpenalty=10000
\displaywidowpenalty=10000
\begin{tabularx}{\linewidth}{@{}>{\bfseries}p{0.22\linewidth}X@{}}
組織 & DeepSeek \\
種別 & - \\
時系列 & 2025-01-20 \\
\end{tabularx}
\smallskip
{\bfseries 一次資料から整理したtechnical points}
\begin{itemize}[leftmargin=1.5em,itemsep=0.35em]
\item \textbf{一次情報で確認できる事実}: 2025年1月20日、DeepSeekはDeepSeek-R1をdeepseek-reasoner endpointとしてAPI提供した。
\end{itemize}
\begin{samepage}
{\bfseries 一次資料}
\begingroup\sloppy
\Urlmuskip=0mu plus 2mu\relax
\begin{itemize}[leftmargin=1.5em,itemsep=0.25em]
\item {\scriptsize\url{https://api-docs.deepseek.com/updates}}
\end{itemize}
\endgroup
\end{samepage}
\endgroup
\end{technicalnote}
\medskip

\begin{technicalnote}{OpenAI Operator research preview}{補足資料}
\begingroup
% reader-facing Technical Notes break policy
\widowpenalty=10000
\clubpenalty=10000
\displaywidowpenalty=10000
\begin{tabularx}{\linewidth}{@{}>{\bfseries}p{0.22\linewidth}X@{}}
組織 & OpenAI \\
種別 & - \\
時系列 & 2025-01-23 \\
\end{tabularx}
\smallskip
{\bfseries 一次資料から整理したtechnical points}
\begin{itemize}[leftmargin=1.5em,itemsep=0.35em]
\item \textbf{一次情報で確認できる事実}: 2025年1月23日、OpenAIは米国のProユーザー向けにブラウザ操作Agent「Operator」のresearch previewを開始した。
\end{itemize}
\begin{samepage}
{\bfseries 一次資料}
\begingroup\sloppy
\Urlmuskip=0mu plus 2mu\relax
\begin{itemize}[leftmargin=1.5em,itemsep=0.25em]
\item {\scriptsize\url{https://openai.com/index/introducing-operator/}}
\end{itemize}
\endgroup
\end{samepage}
\endgroup
\end{technicalnote}
\medskip

\begin{technicalnote}{Qwen2.5-VL open models}{補足資料}
\begingroup
% reader-facing Technical Notes break policy
\widowpenalty=10000
\clubpenalty=10000
\displaywidowpenalty=10000
\begin{tabularx}{\linewidth}{@{}>{\bfseries}p{0.22\linewidth}X@{}}
組織 & Qwen \\
種別 & - \\
時系列 & 2025-01-26 \\
\end{tabularx}
\smallskip
{\bfseries 一次資料から整理したtechnical points}
\begin{itemize}[leftmargin=1.5em,itemsep=0.35em]
\item \textbf{一次情報で確認できる事実}: 2025年1月26日、Qwenは3B・7B・72BのQwen2.5-VL vision-languageモデルを公開した。
\end{itemize}
\begin{samepage}
{\bfseries 一次資料}
\begingroup\sloppy
\Urlmuskip=0mu plus 2mu\relax
\begin{itemize}[leftmargin=1.5em,itemsep=0.25em]
\item {\scriptsize\url{https://qwenlm.github.io/blog/qwen2.5-vl/}}
\end{itemize}
\endgroup
\end{samepage}
\endgroup
\end{technicalnote}
\medskip

\begin{technicalnote}{Qwen2.5-Max API availability}{補足資料}
\begingroup
% reader-facing Technical Notes break policy
\widowpenalty=10000
\clubpenalty=10000
\displaywidowpenalty=10000
\begin{tabularx}{\linewidth}{@{}>{\bfseries}p{0.22\linewidth}X@{}}
組織 & Qwen \\
種別 & - \\
時系列 & 2025-01-28 \\
\end{tabularx}
\smallskip
{\bfseries 一次資料から整理したtechnical points}
\begin{itemize}[leftmargin=1.5em,itemsep=0.35em]
\item \textbf{一次情報で確認できる事実}: 2025年1月28日、Qwenは大規模MoEモデルQwen2.5-Maxを発表し、Alibaba Cloud APIで提供した。
\end{itemize}
\begin{samepage}
{\bfseries 一次資料}
\begingroup\sloppy
\Urlmuskip=0mu plus 2mu\relax
\begin{itemize}[leftmargin=1.5em,itemsep=0.25em]
\item {\scriptsize\url{https://qwenlm.github.io/blog/qwen2.5-max/}}
\end{itemize}
\endgroup
\end{samepage}
\endgroup
\end{technicalnote}
\medskip

\begin{technicalnote}{OpenAI o3-mini}{補足資料}
\begingroup
% reader-facing Technical Notes break policy
\widowpenalty=10000
\clubpenalty=10000
\displaywidowpenalty=10000
\begin{tabularx}{\linewidth}{@{}>{\bfseries}p{0.22\linewidth}X@{}}
組織 & OpenAI \\
種別 & - \\
時系列 & 2025-01-31 \\
\end{tabularx}
\smallskip
{\bfseries 一次資料から整理したtechnical points}
\begin{itemize}[leftmargin=1.5em,itemsep=0.35em]
\item \textbf{一次情報で確認できる事実}: 2025年1月31日、OpenAIは小型reasoningモデルo3-miniをChatGPTとAPIで提供した。
\end{itemize}
\begin{samepage}
{\bfseries 一次資料}
\begingroup\sloppy
\Urlmuskip=0mu plus 2mu\relax
\begin{itemize}[leftmargin=1.5em,itemsep=0.25em]
\item {\scriptsize\url{https://openai.com/index/openai-o3-mini/}}
\end{itemize}
\endgroup
\end{samepage}
\endgroup
\end{technicalnote}
\medskip

\begin{technicalnote}{OpenAI deep research}{補足資料}
\begingroup
% reader-facing Technical Notes break policy
\widowpenalty=10000
\clubpenalty=10000
\displaywidowpenalty=10000
\begin{tabularx}{\linewidth}{@{}>{\bfseries}p{0.22\linewidth}X@{}}
組織 & OpenAI \\
種別 & - \\
時系列 & 2025-02-02 \\
\end{tabularx}
\smallskip
{\bfseries 一次資料から整理したtechnical points}
\begin{itemize}[leftmargin=1.5em,itemsep=0.35em]
\item \textbf{一次情報で確認できる事実}: 2025年2月2日、OpenAIはChatGPTで複数段階のWeb調査を行うDeep Researchを公開した。
\end{itemize}
\begin{samepage}
{\bfseries 一次資料}
\begingroup\sloppy
\Urlmuskip=0mu plus 2mu\relax
\begin{itemize}[leftmargin=1.5em,itemsep=0.25em]
\item {\scriptsize\url{https://openai.com/index/introducing-deep-research/}}
\end{itemize}
\endgroup
\end{samepage}
\endgroup
\end{technicalnote}
\medskip

\begin{technicalnote}{Claude 3.7 Sonnet and Claude Code}{補足資料}
\begingroup
% reader-facing Technical Notes break policy
\widowpenalty=10000
\clubpenalty=10000
\displaywidowpenalty=10000
\begin{tabularx}{\linewidth}{@{}>{\bfseries}p{0.22\linewidth}X@{}}
組織 & Anthropic \\
種別 & - \\
時系列 & 2025-02-24 \\
\end{tabularx}
\smallskip
{\bfseries 一次資料から整理したtechnical points}
\begin{itemize}[leftmargin=1.5em,itemsep=0.35em]
\item \textbf{一次情報で確認できる事実}: 2025年2月24日、Anthropicはhybrid reasoningモデルClaude 3.7 Sonnetと、coding AgentのClaude Code research previewを発表した。
\end{itemize}
\begin{samepage}
{\bfseries 一次資料}
\begingroup\sloppy
\Urlmuskip=0mu plus 2mu\relax
\begin{itemize}[leftmargin=1.5em,itemsep=0.25em]
\item {\scriptsize\url{https://www.anthropic.com/news/claude-3-7-sonnet}}
\end{itemize}
\endgroup
\end{samepage}
\endgroup
\end{technicalnote}
\medskip

\begin{technicalnote}{GPT-4.5 research preview}{補足資料}
\begingroup
% reader-facing Technical Notes break policy
\widowpenalty=10000
\clubpenalty=10000
\displaywidowpenalty=10000
\begin{tabularx}{\linewidth}{@{}>{\bfseries}p{0.22\linewidth}X@{}}
組織 & OpenAI \\
種別 & - \\
時系列 & 2025-02-27 \\
\end{tabularx}
\smallskip
{\bfseries 一次資料から整理したtechnical points}
\begin{itemize}[leftmargin=1.5em,itemsep=0.35em]
\item \textbf{一次情報で確認できる事実}: 2025年2月27日、OpenAIはGPT-4.5をresearch previewとして公開した。
\end{itemize}
\begin{samepage}
{\bfseries 一次資料}
\begingroup\sloppy
\Urlmuskip=0mu plus 2mu\relax
\begin{itemize}[leftmargin=1.5em,itemsep=0.25em]
\item {\scriptsize\url{https://openai.com/index/introducing-gpt-4-5/}}
\end{itemize}
\endgroup
\end{samepage}
\endgroup
\end{technicalnote}
\medskip

\begin{technicalnote}{Responses API, built-in tools, and Agents SDK}{補足資料}
\begingroup
% reader-facing Technical Notes break policy
\widowpenalty=10000
\clubpenalty=10000
\displaywidowpenalty=10000
\begin{tabularx}{\linewidth}{@{}>{\bfseries}p{0.22\linewidth}X@{}}
組織 & OpenAI \\
種別 & - \\
時系列 & 2025-03-11 \\
\end{tabularx}
\smallskip
{\bfseries 一次資料から整理したtechnical points}
\begin{itemize}[leftmargin=1.5em,itemsep=0.35em]
\item \textbf{一次情報で確認できる事実}: 2025年3月11日、OpenAIはResponses API、Web/File Search、Computer Use、Agents SDK、TracingなどのAgent構築機能を発表した。
\end{itemize}
\begin{samepage}
{\bfseries 一次資料}
\begingroup\sloppy
\Urlmuskip=0mu plus 2mu\relax
\begin{itemize}[leftmargin=1.5em,itemsep=0.25em]
\item {\scriptsize\url{https://openai.com/index/new-tools-for-building-agents/}}
\end{itemize}
\endgroup
\end{samepage}
\endgroup
\end{technicalnote}
\medskip

\begin{technicalnote}{Gemma 3 open models}{補足資料}
\begingroup
% reader-facing Technical Notes break policy
\widowpenalty=10000
\clubpenalty=10000
\displaywidowpenalty=10000
\begin{tabularx}{\linewidth}{@{}>{\bfseries}p{0.22\linewidth}X@{}}
組織 & Google \\
種別 & - \\
時系列 & 2025-03-12 \\
\end{tabularx}
\smallskip
{\bfseries 一次資料から整理したtechnical points}
\begin{itemize}[leftmargin=1.5em,itemsep=0.35em]
\item \textbf{一次情報で確認できる事実}: 2025年3月12日、Googleは1B・4B・12B・27BのGemma 3 open model familyを公開した。
\end{itemize}
\begin{samepage}
{\bfseries 一次資料}
\begingroup\sloppy
\Urlmuskip=0mu plus 2mu\relax
\begin{itemize}[leftmargin=1.5em,itemsep=0.25em]
\item {\scriptsize\url{https://blog.google/innovation-and-ai/technology/developers-tools/gemma-3/}}
\end{itemize}
\endgroup
\end{samepage}
\endgroup
\end{technicalnote}
\medskip

\begin{technicalnote}{Claude web search product preview}{補足資料}
\begingroup
% reader-facing Technical Notes break policy
\widowpenalty=10000
\clubpenalty=10000
\displaywidowpenalty=10000
\begin{tabularx}{\linewidth}{@{}>{\bfseries}p{0.22\linewidth}X@{}}
組織 & Anthropic \\
種別 & - \\
時系列 & 2025-03-20 \\
\end{tabularx}
\smallskip
{\bfseries 一次資料から整理したtechnical points}
\begin{itemize}[leftmargin=1.5em,itemsep=0.35em]
\item \textbf{一次情報で確認できる事実}: 2025年3月20日、AnthropicはClaude製品にWeb Searchを追加し、その後5月に提供地域を拡大した。
\end{itemize}
\begin{samepage}
{\bfseries 一次資料}
\begingroup\sloppy
\Urlmuskip=0mu plus 2mu\relax
\begin{itemize}[leftmargin=1.5em,itemsep=0.25em]
\item {\scriptsize\url{https://www.anthropic.com/news/web-search}}
\end{itemize}
\endgroup
\end{samepage}
\endgroup
\end{technicalnote}
\medskip

\begin{technicalnote}{DeepSeek-V3-0324 API upgrade}{補足資料}
\begingroup
% reader-facing Technical Notes break policy
\widowpenalty=10000
\clubpenalty=10000
\displaywidowpenalty=10000
\begin{tabularx}{\linewidth}{@{}>{\bfseries}p{0.22\linewidth}X@{}}
組織 & DeepSeek \\
種別 & - \\
時系列 & 2025-03-24 \\
\end{tabularx}
\smallskip
{\bfseries 一次資料から整理したtechnical points}
\begin{itemize}[leftmargin=1.5em,itemsep=0.35em]
\item \textbf{一次情報で確認できる事実}: 2025年3月24日、DeepSeekはdeepseek-chat serviceをV3-0324へ更新した。
\end{itemize}
\begin{samepage}
{\bfseries 一次資料}
\begingroup\sloppy
\Urlmuskip=0mu plus 2mu\relax
\begin{itemize}[leftmargin=1.5em,itemsep=0.25em]
\item {\scriptsize\url{https://api-docs.deepseek.com/updates}}
\end{itemize}
\endgroup
\end{samepage}
\endgroup
\end{technicalnote}
\medskip

\begin{technicalnote}{Gemini 2.5 Pro Experimental}{補足資料}
\begingroup
% reader-facing Technical Notes break policy
\widowpenalty=10000
\clubpenalty=10000
\displaywidowpenalty=10000
\begin{tabularx}{\linewidth}{@{}>{\bfseries}p{0.22\linewidth}X@{}}
組織 & Google \\
種別 & - \\
時系列 & 2025-03-25 \\
\end{tabularx}
\smallskip
{\bfseries 一次資料から整理したtechnical points}
\begin{itemize}[leftmargin=1.5em,itemsep=0.35em]
\item \textbf{一次情報で確認できる事実}: 2025年3月25日、Googleはthinking、long context、coding、multimodal capabilityを持つGemini 2.5 Pro Experimentalを公開した。
\end{itemize}
\begin{samepage}
{\bfseries 一次資料}
\begingroup\sloppy
\Urlmuskip=0mu plus 2mu\relax
\begin{itemize}[leftmargin=1.5em,itemsep=0.25em]
\item {\scriptsize\url{https://blog.google/innovation-and-ai/models-and-research/google-deepmind/gemini-model-thinking-updates-march-2025/}}
\end{itemize}
\endgroup
\end{samepage}
\endgroup
\end{technicalnote}
\medskip

\begin{technicalnote}{GPT-4o native image generation}{補足資料}
\begingroup
% reader-facing Technical Notes break policy
\widowpenalty=10000
\clubpenalty=10000
\displaywidowpenalty=10000
\begin{tabularx}{\linewidth}{@{}>{\bfseries}p{0.22\linewidth}X@{}}
組織 & OpenAI \\
種別 & - \\
時系列 & 2025-03-25 \\
\end{tabularx}
\smallskip
{\bfseries 一次資料から整理したtechnical points}
\begin{itemize}[leftmargin=1.5em,itemsep=0.35em]
\item \textbf{一次情報で確認できる事実}: 2025年3月25日、OpenAIはGPT-4oにnative image generationを統合した。
\end{itemize}
\begin{samepage}
{\bfseries 一次資料}
\begingroup\sloppy
\Urlmuskip=0mu plus 2mu\relax
\begin{itemize}[leftmargin=1.5em,itemsep=0.25em]
\item {\scriptsize\url{https://openai.com/index/introducing-4o-image-generation/}}
\end{itemize}
\endgroup
\end{samepage}
\endgroup
\end{technicalnote}
\medskip

\begin{technicalnote}{Tracing the thoughts of a large language model}{補足資料}
\begingroup
% reader-facing Technical Notes break policy
\widowpenalty=10000
\clubpenalty=10000
\displaywidowpenalty=10000
\begin{tabularx}{\linewidth}{@{}>{\bfseries}p{0.22\linewidth}X@{}}
組織 & Anthropic \\
種別 & - \\
時系列 & 2025-03-27 \\
\end{tabularx}
\smallskip
{\bfseries 一次資料から整理したtechnical points}
\begin{itemize}[leftmargin=1.5em,itemsep=0.35em]
\item \textbf{一次情報で確認できる事実}: 2025年3月27日、Anthropicはlarge language modelの内部思考を追跡するinterpretability研究を公開した。
\end{itemize}
\begin{samepage}
{\bfseries 一次資料}
\begingroup\sloppy
\Urlmuskip=0mu plus 2mu\relax
\begin{itemize}[leftmargin=1.5em,itemsep=0.25em]
\item {\scriptsize\url{https://www.anthropic.com/research/tracing-thoughts-language-model}}
\end{itemize}
\endgroup
\end{samepage}
\endgroup
\end{technicalnote}
\medskip

\begin{technicalnote}{Qwen2.5-Omni}{補足資料}
\begingroup
% reader-facing Technical Notes break policy
\widowpenalty=10000
\clubpenalty=10000
\displaywidowpenalty=10000
\begin{tabularx}{\linewidth}{@{}>{\bfseries}p{0.22\linewidth}X@{}}
組織 & Qwen \\
種別 & - \\
時系列 & 2025-03-27 \\
\end{tabularx}
\smallskip
{\bfseries 一次資料から整理したtechnical points}
\begin{itemize}[leftmargin=1.5em,itemsep=0.35em]
\item \textbf{一次情報で確認できる事実}: 2025年3月27日、Qwenはtext・image・audio・video入力とstreaming speech出力を扱うQwen2.5-Omniを公開した。
\end{itemize}
\begin{samepage}
{\bfseries 一次資料}
\begingroup\sloppy
\Urlmuskip=0mu plus 2mu\relax
\begin{itemize}[leftmargin=1.5em,itemsep=0.25em]
\item {\scriptsize\url{https://qwenlm.github.io/blog/qwen2.5-omni/}}
\end{itemize}
\endgroup
\end{samepage}
\endgroup
\end{technicalnote}
\medskip

\begin{technicalnote}{QVQ-Max visual reasoning}{補足資料}
\begingroup
% reader-facing Technical Notes break policy
\widowpenalty=10000
\clubpenalty=10000
\displaywidowpenalty=10000
\begin{tabularx}{\linewidth}{@{}>{\bfseries}p{0.22\linewidth}X@{}}
組織 & Qwen \\
種別 & - \\
時系列 & 2025-03-28 \\
\end{tabularx}
\smallskip
{\bfseries 一次資料から整理したtechnical points}
\begin{itemize}[leftmargin=1.5em,itemsep=0.35em]
\item \textbf{一次情報で確認できる事実}: 2025年3月28日、Qwenはvisual reasoningモデルQVQ-Maxを公開した。
\end{itemize}
\begin{samepage}
{\bfseries 一次資料}
\begingroup\sloppy
\Urlmuskip=0mu plus 2mu\relax
\begin{itemize}[leftmargin=1.5em,itemsep=0.25em]
\item {\scriptsize\url{https://qwenlm.github.io/blog/qvq-max/}}
\end{itemize}
\endgroup
\end{samepage}
\endgroup
\end{technicalnote}
\medskip

\begin{technicalnote}{Llama 4 Scout and Maverick}{補足資料}
\begingroup
% reader-facing Technical Notes break policy
\widowpenalty=10000
\clubpenalty=10000
\displaywidowpenalty=10000
\begin{tabularx}{\linewidth}{@{}>{\bfseries}p{0.22\linewidth}X@{}}
組織 & Meta \\
種別 & - \\
時系列 & 2025-04-05 \\
\end{tabularx}
\smallskip
{\bfseries 一次資料から整理したtechnical points}
\begin{itemize}[leftmargin=1.5em,itemsep=0.35em]
\item \textbf{一次情報で確認できる事実}: 2025年4月5日、Metaはopen-weightのnatively multimodal MoEモデルLlama 4 ScoutとMaverickを公開した。
\end{itemize}
\begin{samepage}
{\bfseries 一次資料}
\begingroup\sloppy
\Urlmuskip=0mu plus 2mu\relax
\begin{itemize}[leftmargin=1.5em,itemsep=0.25em]
\item {\scriptsize\url{https://ai.meta.com/blog/llama-4-multimodal-intelligence/}}
\end{itemize}
\endgroup
\end{samepage}
\endgroup
\end{technicalnote}
\medskip

\begin{technicalnote}{GPT-4.1 API family}{補足資料}
\begingroup
% reader-facing Technical Notes break policy
\widowpenalty=10000
\clubpenalty=10000
\displaywidowpenalty=10000
\begin{tabularx}{\linewidth}{@{}>{\bfseries}p{0.22\linewidth}X@{}}
組織 & OpenAI \\
種別 & - \\
時系列 & 2025-04-14 \\
\end{tabularx}
\smallskip
{\bfseries 一次資料から整理したtechnical points}
\begin{itemize}[leftmargin=1.5em,itemsep=0.35em]
\item \textbf{一次情報で確認できる事実}: 2025年4月14日、OpenAIはGPT-4.1、mini、nanoのAPI familyを公開した。
\end{itemize}
\begin{samepage}
{\bfseries 一次資料}
\begingroup\sloppy
\Urlmuskip=0mu plus 2mu\relax
\begin{itemize}[leftmargin=1.5em,itemsep=0.25em]
\item {\scriptsize\url{https://openai.com/index/gpt-4-1/}}
\end{itemize}
\endgroup
\end{samepage}
\endgroup
\end{technicalnote}
\medskip

\begin{technicalnote}{GPT-4.5 API deprecation announced}{補足資料}
\begingroup
% reader-facing Technical Notes break policy
\widowpenalty=10000
\clubpenalty=10000
\displaywidowpenalty=10000
\begin{tabularx}{\linewidth}{@{}>{\bfseries}p{0.22\linewidth}X@{}}
組織 & OpenAI \\
種別 & - \\
時系列 & 2025-04-14 \\
\end{tabularx}
\smallskip
{\bfseries 一次資料から整理したtechnical points}
\begin{itemize}[leftmargin=1.5em,itemsep=0.35em]
\item \textbf{一次情報で確認できる事実}: 2025年4月14日、OpenAIはGPT-4.5 Preview APIを7月14日にretireする予定を告知した。
\end{itemize}
\begin{samepage}
{\bfseries 一次資料}
\begingroup\sloppy
\Urlmuskip=0mu plus 2mu\relax
\begin{itemize}[leftmargin=1.5em,itemsep=0.25em]
\item {\scriptsize\url{https://openai.com/index/gpt-4-1/}}
\end{itemize}
\endgroup
\end{samepage}
\endgroup
\end{technicalnote}
\medskip

\begin{technicalnote}{OpenAI o3 and o4-mini}{補足資料}
\begingroup
% reader-facing Technical Notes break policy
\widowpenalty=10000
\clubpenalty=10000
\displaywidowpenalty=10000
\begin{tabularx}{\linewidth}{@{}>{\bfseries}p{0.22\linewidth}X@{}}
組織 & OpenAI \\
種別 & - \\
時系列 & 2025-04-16 \\
\end{tabularx}
\smallskip
{\bfseries 一次資料から整理したtechnical points}
\begin{itemize}[leftmargin=1.5em,itemsep=0.35em]
\item \textbf{一次情報で確認できる事実}: 2025年4月16日、OpenAIはtool accessを備えたreasoningモデルo3とo4-miniを公開し、同日にCodex CLIも発表した。
\end{itemize}
\begin{samepage}
{\bfseries 一次資料}
\begingroup\sloppy
\Urlmuskip=0mu plus 2mu\relax
\begin{itemize}[leftmargin=1.5em,itemsep=0.25em]
\item {\scriptsize\url{https://openai.com/index/introducing-o3-and-o4-mini/}}
\end{itemize}
\endgroup
\end{samepage}
\endgroup
\end{technicalnote}
\medskip

\begin{technicalnote}{Llama API limited preview}{補足資料}
\begingroup
% reader-facing Technical Notes break policy
\widowpenalty=10000
\clubpenalty=10000
\displaywidowpenalty=10000
\begin{tabularx}{\linewidth}{@{}>{\bfseries}p{0.22\linewidth}X@{}}
組織 & Meta \\
種別 & - \\
時系列 & 2025-04-29 \\
\end{tabularx}
\smallskip
{\bfseries 一次資料から整理したtechnical points}
\begin{itemize}[leftmargin=1.5em,itemsep=0.35em]
\item \textbf{一次情報で確認できる事実}: 2025年4月29日、MetaはLlama ecosystem向けhosted API/toolingのlimited previewを開始した。
\end{itemize}
\begin{samepage}
{\bfseries 一次資料}
\begingroup\sloppy
\Urlmuskip=0mu plus 2mu\relax
\begin{itemize}[leftmargin=1.5em,itemsep=0.25em]
\item {\scriptsize\url{https://ai.meta.com/blog/llamacon-llama-news/}}
\end{itemize}
\endgroup
\end{samepage}
\endgroup
\end{technicalnote}
\medskip

\begin{technicalnote}{Llama Guard 4 and protection tools}{補足資料}
\begingroup
% reader-facing Technical Notes break policy
\widowpenalty=10000
\clubpenalty=10000
\displaywidowpenalty=10000
\begin{tabularx}{\linewidth}{@{}>{\bfseries}p{0.22\linewidth}X@{}}
組織 & Meta \\
種別 & - \\
時系列 & 2025-04-29 \\
\end{tabularx}
\smallskip
{\bfseries 一次資料から整理したtechnical points}
\begin{itemize}[leftmargin=1.5em,itemsep=0.35em]
\item \textbf{一次情報で確認できる事実}: 2025年4月29日、Metaはmultimodal safeguardモデルLlama Guard 4とprotection toolsを公開した。
\end{itemize}
\begin{samepage}
{\bfseries 一次資料}
\begingroup\sloppy
\Urlmuskip=0mu plus 2mu\relax
\begin{itemize}[leftmargin=1.5em,itemsep=0.25em]
\item {\scriptsize\url{https://ai.meta.com/blog/ai-defenders-program-llama-protection-tools/}}
\end{itemize}
\endgroup
\end{samepage}
\endgroup
\end{technicalnote}
\medskip

\begin{technicalnote}{Qwen3}{補足資料}
\begingroup
% reader-facing Technical Notes break policy
\widowpenalty=10000
\clubpenalty=10000
\displaywidowpenalty=10000
\begin{tabularx}{\linewidth}{@{}>{\bfseries}p{0.22\linewidth}X@{}}
組織 & Qwen \\
種別 & - \\
時系列 & 2025-04-29 \\
\end{tabularx}
\smallskip
{\bfseries 一次資料から整理したtechnical points}
\begin{itemize}[leftmargin=1.5em,itemsep=0.35em]
\item \textbf{一次情報で確認できる事実}: 2025年4月29日、Qwenはthinking/non-thinking modeを持つopen-weight dense/MoE family Qwen3を公開した。
\end{itemize}
\begin{samepage}
{\bfseries 一次資料}
\begingroup\sloppy
\Urlmuskip=0mu plus 2mu\relax
\begin{itemize}[leftmargin=1.5em,itemsep=0.25em]
\item {\scriptsize\url{https://qwenlm.github.io/blog/qwen3/}}
\end{itemize}
\endgroup
\end{samepage}
\endgroup
\end{technicalnote}
\medskip

\begin{technicalnote}{Anthropic web search API}{補足資料}
\begingroup
% reader-facing Technical Notes break policy
\widowpenalty=10000
\clubpenalty=10000
\displaywidowpenalty=10000
\begin{tabularx}{\linewidth}{@{}>{\bfseries}p{0.22\linewidth}X@{}}
組織 & Anthropic \\
種別 & - \\
時系列 & 2025-05-07 \\
\end{tabularx}
\smallskip
{\bfseries 一次資料から整理したtechnical points}
\begin{itemize}[leftmargin=1.5em,itemsep=0.35em]
\item \textbf{一次情報で確認できる事実}: 2025年5月7日、AnthropicはMessages APIにWeb Search toolを追加した。
\end{itemize}
\begin{samepage}
{\bfseries 一次資料}
\begingroup\sloppy
\Urlmuskip=0mu plus 2mu\relax
\begin{itemize}[leftmargin=1.5em,itemsep=0.25em]
\item {\scriptsize\url{https://www.anthropic.com/news/web-search-api}}
\end{itemize}
\endgroup
\end{samepage}
\endgroup
\end{technicalnote}
\medskip

\begin{technicalnote}{OpenAI Codex cloud software engineering agent}{補足資料}
\begingroup
% reader-facing Technical Notes break policy
\widowpenalty=10000
\clubpenalty=10000
\displaywidowpenalty=10000
\begin{tabularx}{\linewidth}{@{}>{\bfseries}p{0.22\linewidth}X@{}}
組織 & OpenAI \\
種別 & - \\
時系列 & 2025-05-16 \\
\end{tabularx}
\smallskip
{\bfseries 一次資料から整理したtechnical points}
\begin{itemize}[leftmargin=1.5em,itemsep=0.35em]
\item \textbf{一次情報で確認できる事実}: 2025年5月16日、OpenAIは並列sandboxでタスクを実行するcloud software-engineering Agent Codexをresearch previewとして公開した。
\end{itemize}
\begin{samepage}
{\bfseries 一次資料}
\begingroup\sloppy
\Urlmuskip=0mu plus 2mu\relax
\begin{itemize}[leftmargin=1.5em,itemsep=0.25em]
\item {\scriptsize\url{https://openai.com/index/introducing-codex/}}
\end{itemize}
\endgroup
\end{samepage}
\endgroup
\end{technicalnote}
\medskip

\begin{technicalnote}{Jules coding agent public beta}{補足資料}
\begingroup
% reader-facing Technical Notes break policy
\widowpenalty=10000
\clubpenalty=10000
\displaywidowpenalty=10000
\begin{tabularx}{\linewidth}{@{}>{\bfseries}p{0.22\linewidth}X@{}}
組織 & Google \\
種別 & - \\
時系列 & 2025-05-20 \\
\end{tabularx}
\smallskip
{\bfseries 一次資料から整理したtechnical points}
\begin{itemize}[leftmargin=1.5em,itemsep=0.35em]
\item \textbf{一次情報で確認できる事実}: 2025年5月20日、GoogleはGitHub repositoryをcloud VM上で扱う非同期coding Agent Julesをpublic betaで公開した。
\end{itemize}
\begin{samepage}
{\bfseries 一次資料}
\begingroup\sloppy
\Urlmuskip=0mu plus 2mu\relax
\begin{itemize}[leftmargin=1.5em,itemsep=0.25em]
\item {\scriptsize\url{https://developers.googleblog.com/google-io-2025-developer-keynote-recap/}}
\end{itemize}
\endgroup
\end{samepage}
\endgroup
\end{technicalnote}
\medskip

\begin{technicalnote}{Veo 3, Imagen 4, Flow, and Lyria 2}{補足資料}
\begingroup
% reader-facing Technical Notes break policy
\widowpenalty=10000
\clubpenalty=10000
\displaywidowpenalty=10000
\begin{tabularx}{\linewidth}{@{}>{\bfseries}p{0.22\linewidth}X@{}}
組織 & Google \\
種別 & - \\
時系列 & 2025-05-20 \\
\end{tabularx}
\smallskip
{\bfseries 一次資料から整理したtechnical points}
\begin{itemize}[leftmargin=1.5em,itemsep=0.35em]
\item \textbf{一次情報で確認できる事実}: 2025年5月20日、GoogleはVeo 3、Imagen 4、Flow、Lyria 2を含むgenerative media model/tool群を発表した。
\end{itemize}
\begin{samepage}
{\bfseries 一次資料}
\begingroup\sloppy
\Urlmuskip=0mu plus 2mu\relax
\begin{itemize}[leftmargin=1.5em,itemsep=0.25em]
\item {\scriptsize\url{https://blog.google/innovation-and-ai/products/generative-media-models-io-2025/}}
\end{itemize}
\endgroup
\end{samepage}
\endgroup
\end{technicalnote}
\medskip

\begin{technicalnote}{Responses API remote MCP and tool expansion}{補足資料}
\begingroup
% reader-facing Technical Notes break policy
\widowpenalty=10000
\clubpenalty=10000
\displaywidowpenalty=10000
\begin{tabularx}{\linewidth}{@{}>{\bfseries}p{0.22\linewidth}X@{}}
組織 & OpenAI \\
種別 & - \\
時系列 & 2025-05-21 \\
\end{tabularx}
\smallskip
{\bfseries 一次資料から整理したtechnical points}
\begin{itemize}[leftmargin=1.5em,itemsep=0.35em]
\item \textbf{一次情報で確認できる事実}: 2025年5月21日、OpenAIはResponses APIにremote MCP、image generation、Code Interpreter、background mode、reasoning summariesを追加した。
\end{itemize}
\begin{samepage}
{\bfseries 一次資料}
\begingroup\sloppy
\Urlmuskip=0mu plus 2mu\relax
\begin{itemize}[leftmargin=1.5em,itemsep=0.25em]
\item {\scriptsize\url{https://openai.com/index/new-tools-and-features-in-the-responses-api/}}
\end{itemize}
\endgroup
\end{samepage}
\endgroup
\end{technicalnote}
\medskip

\begin{technicalnote}{Claude Opus 4 and Sonnet 4}{補足資料}
\begingroup
% reader-facing Technical Notes break policy
\widowpenalty=10000
\clubpenalty=10000
\displaywidowpenalty=10000
\begin{tabularx}{\linewidth}{@{}>{\bfseries}p{0.22\linewidth}X@{}}
組織 & Anthropic \\
種別 & - \\
時系列 & 2025-05-22 \\
\end{tabularx}
\smallskip
{\bfseries 一次資料から整理したtechnical points}
\begin{itemize}[leftmargin=1.5em,itemsep=0.35em]
\item \textbf{一次情報で確認できる事実}: 2025年5月22日、Anthropicはcoding、reasoning、Agent、tool useを重視したClaude Opus 4とSonnet 4を公開した。
\end{itemize}
\begin{samepage}
{\bfseries 一次資料}
\begingroup\sloppy
\Urlmuskip=0mu plus 2mu\relax
\begin{itemize}[leftmargin=1.5em,itemsep=0.25em]
\item {\scriptsize\url{https://www.anthropic.com/news/claude-4}}
\end{itemize}
\endgroup
\end{samepage}
\endgroup
\end{technicalnote}
\medskip

\begin{technicalnote}{Gemini 2.5 Pro and Flash generally available}{補足資料}
\begingroup
% reader-facing Technical Notes break policy
\widowpenalty=10000
\clubpenalty=10000
\displaywidowpenalty=10000
\begin{tabularx}{\linewidth}{@{}>{\bfseries}p{0.22\linewidth}X@{}}
組織 & Google \\
種別 & - \\
時系列 & 2025-06-17 \\
\end{tabularx}
\smallskip
{\bfseries 一次資料から整理したtechnical points}
\begin{itemize}[leftmargin=1.5em,itemsep=0.35em]
\item \textbf{一次情報で確認できる事実}: 2025年6月17日、GoogleはGemini 2.5 ProとFlashをGAとし、Flash-Liteをpreviewで公開した。
\end{itemize}
\begin{samepage}
{\bfseries 一次資料}
\begingroup\sloppy
\Urlmuskip=0mu plus 2mu\relax
\begin{itemize}[leftmargin=1.5em,itemsep=0.25em]
\item {\scriptsize\url{https://developers.googleblog.com/en/gemini-2-5-thinking-model-updates/}}
\end{itemize}
\endgroup
\end{samepage}
\endgroup
\end{technicalnote}
\medskip
